\usepackage[utf8]{inputenc}

\usepackage{geometry}
%%%% ARXIV GEOMETRY
\geometry{
    a4paper,
    right=35mm,
    left=35mm,
    top=50mm,
    bottom=55mm,
    footskip=15mm,
}
\renewcommand{\baselinestretch}{1.5} % required by Imperial

\setlength{\parindent}{0em}
\setlength{\parskip}{0em}

\usepackage{float}
\usepackage[justification=centering]{caption}
\usepackage{xcolor}
% colours selected from https://www.rapidtables.com/web/color/RGB_Color.html
\definecolor{main}{RGB}{0,25,51}   % text colour
\definecolor{alt3}{RGB}{0,128,255} % all links
\color{main} % set to main colour

\usepackage[pdftex,
    bookmarks=true,         % show bookmarks bar?
    unicode=true,          % non-Latin characters in Acrobat’s bookmarks
    pdftoolbar=false,        % show Acrobat’s toolbar?
    pdfmenubar=false,        % show Acrobat’s menu?
    pdffitwindow=true,     % window fit to page when opened
    pdfstartview={FitV},    % fits the height of the page to the window
    pdftitle={On spatially irregular ordinary differential equations and a pathwise volatility modelling framework},    % title
    pdfauthor={Ryan McCrickerd},     % author
    pdfsubject={PhD thesis, Imperial College London, 2021},   % subject of the document
    pdfcreator={},   % creator of the document
    pdfproducer={}, % producer of the document
    pdfkeywords={pathwise volatility, ordinary differential equations, initial value problems, forward uniqueness, random differential equations, functional limit theorems, fractional Brownian motion, Riemann-Liouville, Lévy processes, inverse Gaussian process, normal-inverse Gaussian process, diffusion processes, Cox-Ingersoll-Ross process, option pricing, implied volatility, asset pricing, stochastic volatility, rough volatility, Heston model, Monte Carlo simulation, mathematical finance, quantitative finance, time change}, % list of keywords
    pdfnewwindow=true,      % links in new PDF window
    colorlinks=true,       % false: boxed links; true: colored links
    linktocpage=true,      % link page numbers, not titles, in the contents
    linkcolor=alt3,          % color of internal links (change box color with linkbordercolor)
    citecolor=alt3,        % color of links to bibliography
    filecolor=alt3,      % color of file links
    urlcolor=alt3         % color of external links
]{hyperref}
\usepackage{orcidlink} % ORCID icon beside the author's name

\usepackage{aliascnt}
\usepackage{titlesec}

\titlespacing\section{0pt}{12pt plus 0pt minus 0pt}{12pt plus 0pt minus 0pt}
\titlespacing\subsection{0pt}{8pt plus 0pt minus 0pt}{8pt plus 0pt minus 0pt}
\titlespacing\subsubsection{0pt}{8pt plus 0pt minus 0pt}{8pt plus 0pt minus 0pt}

\titleformat{\section}[hang]{\Large\bfseries\filcenter}{\thesection}{1em}{}
\titleformat{\subsection}[hang]{\large\bfseries\filcenter}{\thesubsection}{1em}{}

\usepackage{fancyhdr}

\usepackage{amsmath,amssymb,amsthm,wasysym,stmaryrd,mathrsfs}
\usepackage[UKenglish]{babel}
% THEOREMING
% Declare theorem-style environments with \mynewtheorem to maintain just one
% counter for numbering, and make them work with \autoref.
\newtheorem{dummy}{***}[section]
\newcommand{\mynewtheorem}[2]{
  \newaliascnt{#1}{dummy}
  \newtheorem{#1}[#1]{#2}
  \aliascntresetthe{#1}
  % make \autoref use this name, overriding any default
  \expandafter\def\csname #1autorefname\endcsname{#2}
}

\newtheoremstyle{mythm}%
    {6pt plus 0pt minus 0pt}% Space above
    {2pt plus 0pt minus 0pt}% Space below
    {\itshape}% Body font
    {}% Indent amount
    {\bfseries}% ⟨Theorem head font⟩
    {\bfseries.}% ⟨Punctuation after theorem head ⟩
    {.5em}% Space after theorem head
    {}%

\newtheoremstyle{myrem}%
    {6pt plus 0pt minus 0pt}% Space above
    {2pt plus 0pt minus 0pt}% Space below
    {}% Body font
    {}% Indent amount
    {\bfseries}% ⟨Theorem head font⟩
    {\bfseries.}% ⟨Punctuation after theorem head ⟩
    {.5em}% Space after theorem head
    {}%

\theoremstyle{mythm}
  \mynewtheorem{theorem}{Theorem}
  \mynewtheorem{lemma}{Lemma}
  \mynewtheorem{corollary}{Corollary}

\theoremstyle{myrem}
  \mynewtheorem{example}{Example}
  \mynewtheorem{problem}{Problem}
  \mynewtheorem{definition}{Definition}
  \mynewtheorem{remark}{Remark}

\renewenvironment{proof}
{\textbf{Proof.}\begin{mdseries}}{\qed\end{mdseries}}

% Appear with capital
\def\sectionautorefname{Chapter}
\def\subsectionautorefname{Section}
\def\subsubsectionautorefname{Section}

\usepackage[longnamesfirst]{natbib}

\usepackage{graphicx}

\usepackage{bbm} % for indicator function

\usepackage{mathtools}

\usepackage{changepage}
\usepackage[all]{nowidow}
% Keep the thesis's original page layout, which was built with TeX Live 2020.
\makeatletter
% lineno 5.0 and later wraps amsmath's display environments, which lets a page
% break fall just before them. Line numbers are only used in the appendix, where
% there are no displays, so the environments are put back as amsmath defines them.
\def\thesis@displays{equation,equation*,multline,multline*,gather,gather*,align,align*,alignat,alignat*,flalign,flalign*}
\@for\thesis@display:=\thesis@displays\do{%
  \expandafter\let\csname thesis@\thesis@display\expandafter\endcsname\csname\thesis@display\endcsname
  \expandafter\let\csname thesis@end\thesis@display\expandafter\endcsname\csname end\thesis@display\endcsname}
\let\thesis@math@cr@\math@cr@
\let\thesis@mmeasure@\mmeasure@
\usepackage{lineno}
\@for\thesis@display:=\thesis@displays\do{%
  \expandafter\let\csname\thesis@display\expandafter\endcsname\csname thesis@\thesis@display\endcsname
  \expandafter\let\csname end\thesis@display\expandafter\endcsname\csname thesis@end\thesis@display\endcsname}
\let\math@cr@\thesis@math@cr@
\let\mmeasure@\thesis@mmeasure@
% amsmath now steps the equation counter inside the display. In 2020 it did so
% just before, where hyperref's link anchor could take a line of its own above
% the display. This is the 2020 definition.
\renewenvironment{equation}{%
  \incr@eqnum
  \mathdisplay@push
  \st@rredfalse \global\@eqnswtrue
  \mathdisplay{equation}%
}{%
  \endmathdisplay{equation}%
  \mathdisplay@pop
  \ignorespacesafterend
}
\makeatother
\def\linenumberfont{\normalfont\footnotesize\sffamily\color{gray}}

\usepackage{centernot}

\newcommand{\dd}{\mathrm{d}}
\newcommand{\id}{\mathrm{e}}
\newcommand{\RR}{\mathbb{R}}
\newcommand{\II}{\mathbb{I}}
\newcommand{\QQ}{\mathbb{Q}}
\newcommand{\PP}{\mathbb{P}}
\newcommand{\EE}{\mathbb{E}}
\newcommand{\WW}{\mathbb{W}}
\newcommand{\BB}{\mathbb{B}}
\newcommand{\bS}{\mathbb{S}}
\newcommand{\NN}{\mathbb{N}}
\newcommand{\cX}{\mathcal{X}}
\newcommand{\cE}{\mathcal{E}}
\newcommand{\cF}{\mathcal{F}}
\newcommand{\cG}{\mathcal{G}}
\newcommand{\cD}{\mathcal{D}}
\newcommand{\cB}{\mathcal{B}}
\newcommand{\uW}{\mathrm{W}}
\newcommand{\uF}{\mathrm{F}}
\newcommand{\uFe}{\mathrm{F_{\:\!\! e}}}
\newcommand{\uG}{\mathrm{G}}
\newcommand{\uD}{\mathrm{D}}
\newcommand{\uS}{\mathrm{S}}
\newcommand{\uC}{\mathrm{C}}
\newcommand{\uN}{\mathrm{N}}
\newcommand{\uH}{\mathrm{H}}
\newcommand{\uAC}{\mathrm{AC}}
\newcommand{\uE}{\mathrm{E}}
\newcommand{\uM}{\mathrm{M}}
\newcommand{\uL}{\mathrm{L}}
\newcommand{\uJ}{\mathrm{J}}

\newcommand{\tX}{\tilde{X}}

\newcommand{\tW}{\tilde{W}}
\newcommand{\tS}{\tilde{S}}

\newcommand{\cR}{\mathcal{R}}

\DeclareMathOperator*{\Leb}{Leb}

\newcommand{\ve}{\varepsilon}
\newcommand{\vt}{\vartheta}
\newcommand{\vp}{\varphi}
\newcommand{\hvp}{\hat{\varphi}}

\newcommand{\bvp}{\overline{\varphi}}
\newcommand{\ubvp}{\underline{\varphi}}

\newcommand{\bX}{{\overline{X}}}
\newcommand{\bRR}{{\overline{\mathbb{R}}}}

\newcommand{\bPhi}{\overline{\Phi}}

\newcommand{\bXX}{{\overline{X}}}
\newcommand{\tY}{{\tilde{Y}}}

\newcommand{\hY}{{\hat{Y}}}
\newcommand{\hW}{{\hat{W}}}
\newcommand{\hX}{{\hat{X}}}
\newcommand{\hS}{{\hat{S}}}

\newcommand{\cd}{\xrightarrow{\dd}}
\newcommand{\cfd}{\xrightarrow{\mathrm{f.d.}}}
\newcommand{\cp}{\xrightarrow{\mathrm{p}}}
\newcommand{\asc}{\xrightarrow{\mathrm{a.s.}}}
\newcommand{\ed}{\stackrel{\dd}{=}}

\newcommand{\cw}{\xRightarrow{n\to\infty}}

\numberwithin{equation}{section}

\setlength{\abovecaptionskip}{4pt plus 0pt minus 0pt}
\setlength{\belowcaptionskip}{4pt plus 0pt minus 0pt}

% leaves e.g. subsubsections and beyond out of toc
\setcounter{tocdepth}{2}

\newcommand{\ep}{\epsilon}

\newcommand{\sgn}{\mathrm{sgn}}

\newcommand{\tZ}{\tilde{Z}}
\newcommand{\tvt}{\tilde{\vt}}

% this enables
% \begin{changemargin}{0.5cm}{0.5cm}
% \end{changemargin}
% which is how quote works, with this set to 1cm
\def\changemargin#1#2{\list{}{\rightmargin#2\leftmargin#1}\item[]}
\let\endchangemargin=\endlist

\usepackage[T1]{fontenc}
